\documentclass[10pt,a4paper]{amsart}
\usepackage[margin=19mm]{geometry}
\usepackage{amsmath,amssymb,mathtools}
\usepackage[scr=rsfs,cal=euler]{mathalfa}
\usepackage{xfrac}
\usepackage[unicode,hidelinks,bookmarks=false]{hyperref}
\newcommand{\ie}{i.e.\ }
\newcommand{\cf}{cf.\ }
\newcommand{\resp}{respectively\ }
\newcommand{\Subsection}{\subsection}
\providecommand{\nobreakdash}{\nobreak\hbox{-}}
\setlength{\emergencystretch}{2em}
\multlinegap0pt
\usepackage{bm}
\usepackage{bbm}
\usepackage{dsfont}
\usepackage{xspace}
\usepackage{cancel}
\usepackage{mathtools}
\usepackage{amsthm}
\makeatletter
\@ifundefined{theorem}{\newtheorem{theorem}{Theorem}}{}
\@ifundefined{lemma}{\newtheorem{lemma}[theorem]{Lemma}}{}
\@ifundefined{proposition}{\newtheorem{proposition}[theorem]{Proposition}}{}
\makeatother
\newcommand{\R}{\mathbb{R}}
\title[Line, Spiral, Dense in Space]{Line, Spiral, Dense in Space}
\author{Atakora Agoro, Alastair N. Fletcher}
\date{Source: arXiv:2605.28470v1 (CC BY 4.0)}
\begin{document}
\maketitle

\noindent\emph{Source and licence.} Line, Spiral, Dense in Space. Version arXiv:2605.28470v1 (CC BY 4.0), distributed under the Creative Commons Attribution 4.0 International licence, as recorded in the arXiv licence metadata for this exact version and shown on the abstract page https://arxiv.org/abs/2605.28470v1. Full rights notice: licenses/arxiv-2605.28470-CC-BY-4.0.txt.\par\smallskip

\noindent\textit{Editorial scope.} Counted unit: the Lemma whose source label is \texttt{lem:doublepreims} in the article (source0.tex lines 376--381, source label \texttt{lem:doublepreims}), delivered together with the article's own complete original proof of it (source0.tex lines 383--430, one \texttt{proof} environment of 48 lines). No mathematics is written, repaired or re-derived: statement and proof are the article's own text line for line. Required context retained uncounted: lem:4 (lines 230--234); prop:reldist (lines 244--248); lem:radius (lines 340--343); lem:estimate (lines 363--368). Every definition, display and label of those lines is retained unchanged. Omitted from this package and not redistributed: 1--229, 235--243, 249--339, 344--362, 369--375, 382--382, 431--672. Nothing inside a retained excerpt is omitted or altered: every retained body is delivered exactly as the article's lines are written, so no omission receipt of any kind accompanies this package. Citations: the retained excerpts carry no citation command at all, so no bibliography entry of the article is transcribed and no reference is redistributed. Reference adaptation: none was needed and none was made; every reference inside the delivered unit resolves inside it, so no unresolved reference or citation is produced. Printed slips kept literal: no printed word, symbol, hypothesis or number of the counted statement or of its proof was altered. Layout and class: the article's own class is \texttt{amsart} with packages including graphicx, color, fullpage, amsmath, amsthm, amssymb, verbatim, enumerate; the wrapper is the lane's pinned \texttt{amsart} editorial wrapper at 10pt on A4, which restates the theorem-like environments the retained text needs and adds the article's own macro definitions that the retained text needs (\texttt{\textbackslash R}), with extra wrapper packages (bm, bbm, dsfont, xspace, cancel, mathtools). The delivered numbering is the unit's own. Assets: no retained span contains a figure environment, a graphics-inclusion command, a PDF-page inclusion, a table or a TikZ picture; no image file is copied into this package, the package declares no asset dependency and no image file is redistributed with it. Licence: version arXiv:2605.28470v1 of this article is distributed under the Creative Commons Attribution 4.0 International licence, as recorded in the arXiv licence metadata for this exact version and shown on the abstract page https://arxiv.org/abs/2605.28470v1. A full rights notice is delivered at licenses/arxiv-2605.28470-CC-BY-4.0.txt. Characters: every retained excerpt is pure ASCII, so the delivered engine, XeLaTeX with its default Latin Modern text font, reports no missing character.
\begin{lemma}
\label{lem:4}
Let $t_1<t_2$ and let $E$ be the slab $\R^2 \times (t_1,t_2) \subset \R^3$. Then 
\[ D(\mathcal{Z} , E) \leq \lambda(h)^2 e^{t_2-t_1}.\]
\end{lemma}

\begin{proposition}
    \label{prop:reldist}
Suppose that $E\subset \R^n$ is a Lebesgue measurable set with $m(E) > 0$ and $U \subset E$ is a measurable subset. If $f:E \to \R^n$ is a homeomorphism that satisfies $D(f,E) \leq \lambda$, then 
\[ \frac{1}{\lambda^n} \cdot \frac{m(U)}{m(E)} \leq \frac{ m(f(U))}{m(f(E))} \leq \lambda^n \cdot \frac{ m(U)}{m(E)}.  \]
\end{proposition}

\begin{lemma}
    \label{lem:radius}
Any connected component of $\widetilde{V} \subset H$ contains a disk of radius at least $r_0(8\lambda e^{|x_0|})^{-1}$.
\end{lemma}

\begin{lemma}
    \label{lem:estimate}
Given $|x_0| \in \R^+ \setminus  \{ 0 ,1 \}$, there exists $a>0$ such that if $t_1 > \ln \left |\ln |x_0| \right |$ and $t_2 > t_1+a$, then 
\[ \frac{e^{t_2}}{\sqrt{2}} - e^{t_1} > 2\pi .\]
Moreover, we may choose $a$ so that $e^{2a}>3$.
\end{lemma}

\begin{lemma}
\label{lem:doublepreims}
Let $t_1 > \ln \left |\ln |x_0| \right |$ and let $t_2 > t_1+a $, where $a$ is from Lemma \ref{lem:estimate}. Then 
\[ \frac{ m(f^{-1}(U) \cap F(t_1,t_2) ) }{m(F(t_1,t_2))} \geq C,\]
where $C$ depends only on the input data $r_0,|x_0|$ and $\lambda$.
\end{lemma}

\begin{proof}
First, by construction $\mathcal{Z}$ maps $F$ onto a quadrant of the plane $H$. Without loss of generality, we may assume this quadrant is
\[ Q = \{ (x_1,x_2,x_3) : x_1 \geq |x_2|, x_3 = \ln |x_0| \}.\]
Our analysis for the other quadrants runs analogously. Then $\mathcal{Z}$ maps $F(t_1,t_2)$ onto a logarithmic square in $Q$. That is, two sides of $\mathcal{Z} ( F(t_1,t_2))$ run along the lines $x_1 = \pm x_2$, and the other two sides are arcs of circles with radii $e^{t_1}$ and $e^{t_2}$ respectively.

In particular, $\mathcal{Z} ( F(t_1,t_2))$ contains a trapezoid $T$ in $Q$ whose vertical line segments are contained in the lines $x_1 = e^{t_1}$ and $x_1 = e^{t_2}/\sqrt{2}$, and the other boundary line segments are contained in the lines $x_1 = \pm x_2$. By the hypotheses and Lemma \ref{lem:estimate}, it follows that $T$ is non-trivial. 

If we compare the areas of $\mathcal{Z} ( F(t_1,t_2))$ and $T$, we have
\[ m(T) = (e^{t_2}/\sqrt{2} - e^{t_1})(e^{t_2}/\sqrt{2} + e^{t_1}) = \frac{e^{2t_2}}{2} - e^{2t_1} \]
and
\[ m(\mathcal{Z} ( F(t_1,t_2))) = \frac{\pi}{4} (  e^{2t_2} - e^{2t_1}) .\]

It follows that
\[ \frac{  m(\mathcal{Z}( F(t_1,t_2) ) )}{m(T)} = \frac{\pi}{4} \left ( \frac{e^{2(t_2-t_1)} -1}{e^{2(t_t-t_1)}/2 -1} \right ). \]
By the hypotheses, $t_2-t_1\geq a$ and we may assume by Lemma \ref{lem:estimate} that $e^{2a}>3$. The real function $g(t) = \frac{t-1}{t/2-1}$ is decreasing on $(2,\infty)$ and it follows that we obtain
\begin{equation}
    \label{eq:double1}
\frac{  m(\mathcal{Z}( F(t_1,t_2) ) )}{m(T)} \leq \frac{\pi g(3)}{4} < 2\pi.
\end{equation}

Next, Lemma \ref{lem:estimate} implies that the width of the trapezoid $T$ is at least $2\pi$.
Suppose first that $T$ has width between $2\pi$ and $4\pi$ so that 
\begin{equation}
    \label{eq:Tsize}
    m(T) \leq 4\pi(2e^{t_1} + 4\pi).
\end{equation} 
Moreover, as $t_1>\ln | \ln|x_0||$, the condition that $e^{t_2}/\sqrt{2} - e^{t_1} \leq 4\pi$ can be rearranged to
\begin{equation}
    \label{eq:double2}
e^{t_2-t_1} \leq \sqrt{2}\left( 1+\frac{4\pi}{|\ln |x_0||}\right ).
\end{equation}

Then $T$ contains a column made up of squares taken alternately from $G(B_0 \cap H)$ and $G(B_1 \cap H)$. As the left-hand edge of $T$ has length $2e^{t_1}$, it follows that this column has a total of $2\lceil e^{t_1} \rceil$ squares. By Lemma \ref{lem:radius}, each pair of squares will contain a disk contained in $\widetilde{V}$ of radius at least $r_0(8\lambda e^{|x_0|})^{-1}$.

It follows from this and \eqref{eq:Tsize} that 
\begin{equation}
    \label{eq:Tsmall}
\frac{ m(\widetilde{V}\cap T)}{m(T)} \geq \frac{\lceil e^{t_1} \rceil \cdot \pi r_0^2(8\lambda e^{|x_0|})^{-2} }{ 4\pi ( 2e^{t_1} + 4\pi) } \geq \frac{r_0^2}{2^9 \lambda^2 e^{2|x_0|} }.   
\end{equation}
Putting our various estimates together, by Lemma \ref{lem:4}, Proposition \ref{prop:reldist}, \eqref{eq:double1}, \eqref{eq:double2} and \eqref{eq:Tsmall}, we obtain
\begin{align*}
\frac{ m(f^{-1}(U) \cap F(t_1,t_2) ) }{m(F(t_1,t_2))} &\geq \frac{1}{\lambda^2 e^{t_2-t_1}} \cdot \frac{ m(\widetilde{V} \cap \mathcal{Z}(F(t_1,t_2)) }{m(\mathcal{Z}(F(t_1,t_2)))} \\
&\geq \frac{1}{\lambda^2 e^{t_2-t_1}} \cdot \frac {m(\widetilde{V} \cap T)}{ m(T) \cdot 2\pi} \\
&\geq \frac{ r_0^2}{ 2^{11}\lambda^4 \pi e^{2|x_0|}(1 + 4\pi / |\ln |x_0||) }.
\end{align*}

Now, if $T$ has width larger than $4\pi$, then we can split $T$ up into trapezoids which have width between $2\pi$ and $4\pi$ and so the proportion estimate above will persist.
\end{proof}

\end{document}
